\scanpage{69}\mspage{13}
% archive foliation: 68
avec le morphisme canonique (étale de degré $2$)
\[
f_t:\Delta_t\longrightarrow\underline J
\]
soit l'identité -- i.e. (48) définit une \uline{face} du fibré en octaèdres combinatoires $\Delta_t$ sur $S$. Du coup l'octaèdre est $\underline J$-épinglé, en d'autres termes, il est can. isomorphe au \deleted{$\unsure{I}$} octaèdre standart défini par $\underline J$, \inserted{et une section de $\mu_4^*$} en prenant le revêtement biquadratique $\bar\Sigma_{\underline J,2}$ de $\bar\Sigma_{\underline J}$\ldots
